% !TeX program = xelatex
% ============================================================
%  LukeDoesMaths -- Leaving Certificate notes template (example)
%  Compile with: xelatex notes-template.tex   (or lualatex)
%
%  This file is a worked example showing how to use lukedoesmaths.sty.
%  Copy it as the starting point for a new topic and edit the
%  \subject / \level / \topic fields plus the body content below.
% ============================================================
\documentclass[11pt]{article}
\usepackage{lukedoesmaths}

\subject{Mathematics}
\level{Higher Level}
\topic{Differentiation}

\hypersetup{pdftitle={LukeDoesMaths -- Differentiation Notes},pdfauthor={Luke Johnston}}

\begin{document}

\ldmtitlepage

\tableofcontents
\clearpage

\section{What is a Derivative?}

The derivative of a function tells us its \emph{rate of change} at a
point --- geometrically, the slope of the tangent to the curve at that
point.

\begin{ldmdefinition}
For a function $f(x)$, the derivative $f'(x)$ is defined as
\[
  f'(x) = \lim_{h \to 0} \frac{f(x+h) - f(x)}{h}.
\]
This is sometimes called \emph{differentiation from first principles}
and is examinable at Higher Level.
\end{ldmdefinition}

\begin{ldmexamtip}
First-principles questions almost always appear as a stand-alone
Paper 1 question. Learn the definition \textbf{exactly} as written ---
markers look for the correct limit notation, not just the final answer.
\end{ldmexamtip}

\section{Differentiation Rules}

\begin{ldmformula}
  \begin{tabular}{ll}
    Power rule: & $\dfrac{d}{dx}\left(x^n\right) = nx^{n-1}$ \\[10pt]
    Product rule: & $(uv)' = u'v + uv'$ \\[10pt]
    Quotient rule: & $\left(\dfrac{u}{v}\right)' = \dfrac{u'v - uv'}{v^2}$ \\[10pt]
    Chain rule: & $\dfrac{dy}{dx} = \dfrac{dy}{du}\cdot\dfrac{du}{dx}$
  \end{tabular}
\end{ldmformula}

\begin{ldmexample}
Differentiate $f(x) = 3x^4 - 5x^2 + 2x - 7$ with respect to $x$.

\textbf{Solution.} Apply the power rule to each term:
\[
  f'(x) = 12x^3 - 10x + 2.
\]
\end{ldmexample}

\begin{ldmexample}[Worked Example -- Chain Rule]
Differentiate $g(x) = (3x^2 + 1)^5$.

\textbf{Solution.} Let $u = 3x^2 + 1$, so $g = u^5$.
\[
  \dv{g}{u} = 5u^4, \qquad \dv{u}{x} = 6x
\]
\[
  \dv{g}{x} = 5u^4 \cdot 6x = 30x(3x^2+1)^4.
\]
\end{ldmexample}

\subsection{Common Mistakes}
\begin{itemize}
  \item Forgetting to multiply by the derivative of the inside function (chain rule).
  \item Applying the product rule to a single term multiplied by a constant --- constants just carry through.
  \item Dropping the negative sign when differentiating $\dfrac{1}{x^n} = x^{-n}$.
\end{itemize}

\section{Applications}

\subsection{Applied Maths link: Velocity and Acceleration}

\begin{ldmdefinition}[Definition -- Kinematics]
If $s(t)$ is displacement as a function of time,
\[
  v(t) = \dv{s}{t}, \qquad a(t) = \dv{v}{t} = \dv[2]{s}{t}.
\]
\end{ldmdefinition}

\begin{ldmexample}[Worked Example -- Applied Maths]
A particle moves so that $s(t) = t^3 - 6t^2 + 9t$ (in \unit{\metre}, $t$ in \unit{\second}).
Find the particle's velocity and acceleration at $t = 2\,\unit{\second}$.

\textbf{Solution.}
\[
  v(t) = 3t^2 - 12t + 9 \;\Rightarrow\; v(2) = 12 - 24 + 9 = -3\ \unit{\metre\per\second}
\]
\[
  a(t) = 6t - 12 \;\Rightarrow\; a(2) = 0\ \unit{\metre\per\second\squared}
\]
\end{ldmexample}

\subsection{Physics link: Instantaneous Rate of Change}

\begin{ldmexamtip}
In Physics, differentiation underpins the definitions of instantaneous
velocity and instantaneous acceleration, and appears again in topics
like induced EMF ($\mathcal{E} = -\dv{\Phi}{t}$). Recognising the same
mathematical structure across subjects saves study time.
\end{ldmexamtip}

\section{Summary}

\begin{ldmsummary}
\begin{itemize}
  \item The derivative gives the instantaneous rate of change / gradient of the tangent.
  \item Know the power, product, quotient and chain rules fluently.
  \item First-principles differentiation must use the limit definition exactly.
  \item Differentiation connects directly to kinematics (Applied Maths) and rates of change in Physics.
\end{itemize}
\end{ldmsummary}

\end{document}
